% anexo_a.tex
\chapter{Anexos} 
\label{app:anexo_principal} % Etiqueta para el anexo completo (opcional)

% --- ANEXO A.1: DERIVACIÓN ---
\section{Proof of Proposition \ref{prop:dual_optimal}}
\label{app:derivacion_matematica} % <-- ¡LA ETIQUETA CRÍTICA!

This appendix provides the detailed step-by-step derivation for the dual-network optimization problem, as summarized in Proposition \ref{prop:dual_optimal}.

\subsection{Problem Setup}
The total portfolio $w_i^* = w_i^T + w_i^A$ is exogenous and given. The problem is reformulated as maximizing the agent's utility by choosing the traditional portfolio $w_i^T$, which residually determines the alternative portfolio $w_i^A = w_i^* - w_i^T$.

\subsection{Utility Function}
The total utility function for firm $i$ is the sum of utilities from both networks:
\begin{equation}
g_i(w_i^T, w_i^A, P^T, P^A) = g_i^T(w_i^T, P^T) + g_i^A(w_i^A, P^A)
\end{equation}

Where the utilities in each network are defined as:

\textbf{Traditional Network:}
\begin{equation}
g_i^T(w_i^T, P^T) = (w_i^T)^\top(m_i - P^T e_i) - \gamma_i(w_i^T)^\top\Sigma_i w_i^T
\end{equation}

\textbf{Alternative Network:} (Includes benefit $\tau$ and friction $\lambda_A$)
\begin{equation}
g_i^A(w_i^A, P^A) = (w_i^A)^\top(m_i + \tau - P^A e_i) - \gamma_i(w_i^A)^\top\Sigma_i w_i^A - \lambda_A\|w_i^A\|^2
\end{equation}

\subsection{Optimization Problem}
We substitute $w_i^A = (w_i^* - w_i^T)$ into the total utility function to get a function of $w_i^T$ only:

\begin{align}\label{eq:app_ut_tot}
g_i(w_i^T)
&= (w_i^T)^\top (m_i - P^T e_i)
  - \gamma_i\,(w_i^T)^\top \Sigma_i\,w_i^T \notag\\
&\quad+ (w_i^* - w_i^T)^\top (m_i + \tau - P^A e_i) \notag\\
&\quad- \gamma_i\,(w_i^* - w_i^T)^\top \Sigma_i\,(w_i^* - w_i^T) \\
&\quad - \lambda_A\,\|(w_i^* - w_i^T)\|^2 \notag
\end{align}

This maximization is subject to the traditional network's regulatory constraints:
\begin{equation}
\Psi_i^T w_i^T = w_i^T
\end{equation}

\subsection{Grouping of Terms}
To find the maximum, we first rewrite the objective function $g_i(w_i^T)$ in the standard quadratic form:
\[
g_i(w_i^T)
= \underbrace{\text{constants}}_{C_i}
\;+\;(w_i^T)^\top a_i
\;-\;(w_i^T)^\top H_i\,w_i^T
\]
where \(w := w_i^T\) and \(w^* := w_i^*\).

\subsubsection{1. Expansion of Terms}
We expand the quadratic terms related to $w_i^A = (w_i^* - w_i^T)$:
\begin{align*}
(w^* - w)^\top \Sigma_i\,(w^* - w)
&= (w^*)^\top \Sigma_i\,w^*
  - 2\,(w^*)^\top \Sigma_i\,w
  + w^\top \Sigma_i\,w \\
\|w^* - w\|^2
&= (w^*)^\top w^*
  - 2\,(w^*)^\top w
  + w^\top w
\end{align*}

\subsubsection{2. Identification of Coefficients}
Substituting these expansions into $g_i(w_i^T)$ and grouping terms by $w_i^T$ (as $w$), we identify the constant, linear, and quadratic coefficients.

\begin{itemize}
  \item \textbf{Constant Term ($C_i$):} (All terms not containing $w_i^T$)
    \[
      C_i = (w^*)^\top (m_i + \tau - P^A e_i) 
            - \gamma_i\,(w^*)^\top \Sigma_i\,w^*
            - \lambda_A\,(w^*)^\top w^*.
    \]
  
  \item \textbf{Linear Term ($a_i$):} (All terms linear in $w_i^T$)
    \begin{align*}
      (w_i^T)^\top a_i &= (w_i^T)^\top (m_i - P^T e_i) - (w_i^T)^\top (m_i + \tau - P^A e_i) \\
                      &\quad + 2\gamma_i (w_i^T)^\top \Sigma_i w_i^* + 2\lambda_B (w_i^T)^\top w_i^* \\
      \Rightarrow a_i &= (P^A - P^T)e_i - \tau
                      + 2\,\gamma_i\,\Sigma_i\,w_i^*
                      + 2\,\lambda_A\,w_i^*
    \end{align*}
    
  \item \textbf{Quadratic Term ($H_i$):} (All terms quadratic in $w_i^T$)
    \begin{align*}
      -(w_i^T)^\top H_i w_i^T &= - \gamma_i\,(w_i^T)^\top \Sigma_i\,w_i^T - \gamma_i\,(w_i^T)^\top \Sigma_i\,w_i^T - \lambda_A\,(w_i^T)^\top w_i^T \\
      \Rightarrow H_i &= 2\,\gamma_i\,\Sigma_i + \lambda_A\,I
    \end{align*}
\end{itemize}

\subsubsection{3. Final Quadratic Form}
The objective function is now:
\[
g_i(w_i^T)
= C_i
\;+\;(w_i^T)^\top\,a_i
\;-\;(w_i^T)^\top\,H_i\,w_i^T,
\]
with $a_i$ and $H_i$ defined as in Proposition \ref{prop:dual_optimal} (with $P$ differentiated).

\subsection{Solution with Constraints}
We now solve the constrained problem of maximizing $g_i(w_i^T)$ subject to $\Psi_i^T w_i^T = w_i^T$. This is a standard quadratic optimization problem.

\subsubsection{First-Order Condition}
For the unconstrained maximum, the first-order condition is:
\begin{equation}
\frac{\partial g_i}{\partial w_i^T} = a_i - 2H_i w_i^T = 0 \implies w_i^T = \frac{1}{2}H_i^{-1} a_i
\end{equation}

\subsubsection{Incorporation of Constraints (Lagrangian)}
To incorporate $\Psi_i^T w_i^T = w_i^T$, we use the Lagrangian method:
\begin{equation}
\mathcal{L} = (w_i^T)^\top a_i - (w_i^T)^\top H_i w_i^T - \mu_i^\top(\Psi_i^T w_i^T - w_i^T)
\end{equation}
The first-order conditions are:
\begin{align}
\frac{\partial \mathcal{L}}{\partial w_i^T} &= a_i - 2H_i w_i^T - \Psi_i^{T\top}\mu_i + \mu_i = 0 \\
\frac{\partial \mathcal{L}}{\partial \mu_i} &= \Psi_i^T w_i^T - w_i^T = 0
\end{align}
Solving this system for $w_i^T$ yields the optimal constrained solution.

\subsection{Final Solution}
The solution to the constrained optimization problem is:
\begin{equation}
\boxed{w_i^{T*} = \Psi_i^{T\top}(\Psi_i^T H_i^{-1} \Psi_i^{T\top})^{-1} \Psi_i^T H_i^{-1} a_i}
\end{equation}
where:
\begin{align}
a_i &= (P^A - P^T)e_i - \tau + 2\gamma_i \Sigma_i w_i^* + 2\lambda_A w_i^* \\
H_i &= 2\gamma_i \Sigma_i + \lambda_A I
\end{align}
This expression for $w_i^{T*}$ is exactly what is defined as $Q_i^{\text{dual}} \cdot a_i$ in Proposition \ref{prop:dual_optimal}, completing the proof.

\subsection{Second-Order Condition}
The matrix $H_i = 2\gamma_i \Sigma_i + \lambda_B I$ is positive definite (assuming $\gamma_i > 0$, $\Sigma_i \succ 0$, and $\lambda_B \geq 0$), guaranteeing that the solution found is a unique global maximum.


\section{Example Case with no constraints}
\label{app:3agnoconstraints}
Lets see how changes the network structure when the local firm can access the national bank.

This example presents the baseline scenario where all firms can contract freely without restrictions, allowing us to observe the natural equilibrium structure of the network.

We keep all identical except for the constraint, which now is $\Psi_i=I_n$.

\subsubsection{Pairwise Negotiation Process \label{ex:convergence_nc}}

At the initial state, contracts begin with the optimal values for each firm based on their beliefs, but no payments have been negotiated yet.

\textbf{Initial State ($t = 0$):}

\begin{align}
    W_0 &= \begin{bmatrix} 
        -0.31 & 0.20 & 1.60 \\ 
        0.25 & -0.31 & 1.24 \\ 
        0.33 & 0.44 & -1.94 
    \end{bmatrix} &
    P_0 &= \begin{bmatrix} 
        0 & 0 & 0 \\ 
        0 & 0 & 0 \\ 
        0 & 0 & 0 
    \end{bmatrix}
\end{align}

\textbf{Iteration Process:}

\textbf{Step 5:} After initial negotiations, the network structure begins to stabilize:

\begin{align}
    W_5 &= \begin{bmatrix} 
        -0.75 & -0.06 & 1.04 \\ 
        -0.06 & -0.51 & 0.88 \\ 
        1.02 & 0.87 & -1.31 
    \end{bmatrix} &
    P_5 &= \begin{bmatrix} 
        0.00 & -0.03 & 0.34 \\ 
        0.03 & 0.00 & 0.24 \\ 
        -0.34 & -0.24 & 0.00 
    \end{bmatrix}
\end{align}

The National Bank-Local Firm contract ($W_{13} = 1.04$) emerges as the dominant link. Payment flows have been negotiated with the Local Firm paying both the National Bank ($P_{31} = -0.34$) and the Local Bank ($P_{32} = -0.24$).

\textbf{Step 10:} Approaching equilibrium:

\begin{align}
    W_{10} &= \begin{bmatrix} 
        -0.76 & -0.07 & 1.03 \\ 
        -0.07 & -0.51 & 0.88 \\ 
        1.03 & 0.88 & -1.30 
    \end{bmatrix} &
    P_{10} &= \begin{bmatrix} 
        0.00 & -0.03 & 0.35 \\ 
        0.03 & 0.00 & 0.24 \\ 
        -0.35 & -0.24 & 0.00 
    \end{bmatrix}
\end{align}

Contract sizes have converged to near-final values. The direct connection between National Bank and Local Firm ($W_{13} = 1.03$) remains the largest contract in the network.

\textbf{Convergence:} The algorithm converges to stable values:

\begin{align}
    W_\infty &= \begin{bmatrix} 
        -0.76 & -0.07 & 1.03 \\ 
        -0.07 & -0.51 & 0.88 \\ 
        1.03 & 0.88 & -1.30 
    \end{bmatrix} &
    P_\infty &= \begin{bmatrix} 
        0.00 & -0.03 & 0.35 \\ 
        0.03 & 0.00 & 0.24 \\ 
        -0.35 & -0.24 & 0.00 
    \end{bmatrix}
\end{align}




\subsubsection{No Constraint and 4 agents (input for dual)}
\label{app:optimalw}
\begin{align} 
    W_\infty &= \begin{bmatrix} 
        -0.79 & 0.02 & 1.04 & 0.02 \\ 
        0.02 & -0.17 & 0.56 & -0.17\\ 
        1.03 & 0.56 & -1.45 & 0.56\\
        0.02 & -0.17 & 0.56 & -0.17
    \end{bmatrix} & 
    P_\infty &= \begin{bmatrix} 
        0.00 & 0.19 & 0.44 & 0.19 \\ 
        -0.19 & 0.00 & -0.06 & 0.00 \\ 
        -0.44 & -0.06 & 0.00 & 0.06\\
        -0.19 & 0.00 & -0.06 & 0.00
    \end{bmatrix}
\end{align}